\section{Related Work and Background}
\label{sec:related_work}

Inter-process communication in Unix-like operating systems has been examined extensively from architectural and performance perspectives over several decades~\cite{silberschatz2018operating,kerrisk2010linux}. This section reviews classical IPC taxonomies, analyzes prior performance benchmarks, and highlights the technical gaps addressed by this work.

\subsection{Classical Linux IPC Mechanisms}
The Linux kernel provides two fundamental IPC paradigms: unstructured byte-stream communication and structured message passing~\cite{kerrisk2010linux,stevens1999unixipc}. Anonymous pipes provide unidirectional byte-stream channels managed entirely in memory via kernel page buffers. Full-duplex communication requires two pipes, introducing descriptor management overhead. Named pipes (FIFOs) extend this mechanism to unrelated processes by binding the pipe buffer to a filesystem directory entry created via \texttt{mkfifo()}~\cite{kerrisk2010linux}. Both pipe variants enforce byte-stream semantics, requiring user-space protocols to manage message framing and handle partial read/write conditions.

\subsection{Shared Memory-Based Communication}
POSIX Shared Memory (\texttt{shm\_open}, \texttt{mmap}) represents the highest-throughput IPC mechanism in Unix architectures because it completely eliminates data copying between user and kernel address spaces~\cite{posix2018standard,stevens1999unixipc}. Communicating processes map the same physical memory pages into their respective virtual address spaces, enabling raw memory load/store operations.
However, shared memory lacks implicit synchronization. Processes coordinate access using synchronization primitives such as POSIX unnamed semaphores (\texttt{sem\_init}) backed by the Linux \texttt{futex} subsystem~\cite{franke2002futex}. When uncontended, semaphore acquisition requires only an atomic userspace compare-and-swap instruction; under contention, the kernel suspends the calling thread in a wait queue, introducing scheduler overhead. Furthermore, improper cache-line alignment in shared structures can cause false sharing between CPU cores, degrading cache coherence performance~\cite{herlihy2020art,drepper2007every}.

\subsection{Message Queue Communication}
Structured message queues provide atomic, message-boundary-preserving delivery managed directly by the kernel~\cite{kerrisk2010linux}. 
\textbf{System V Message Queues} (\texttt{msgget}, \texttt{msgsnd}, \texttt{msgrcv}) store messages as linked lists in kernel memory identified by global integer keys, supporting message multiplexing via an integer message type field (\texttt{mtype}). However, System V IPC suffers from a legacy API, lack of file descriptor integration, and absence of native POSIX file permissions~\cite{stevens1999unixipc}.
\textbf{POSIX Message Queues} (\texttt{mq\_open}, \texttt{mq\_send}, \texttt{mq\_receive}) use descriptor semantics (\texttt{mqd\_t}) and are visible within the virtual \texttt{/dev/mqueue} filesystem~\cite{posix2018standard}, providing real-time message priorities (0 to 31) and asynchronous event notification via signals (\texttt{mq\_notify}). Nevertheless, unprivileged Linux configurations enforce strict ceilings on queue depth (defaulting to 10 messages via \texttt{/proc/sys/fs/mqueue/msg\_max}) and require fixed message buffer sizing (\texttt{mq\_msgsize}).

\subsection{Stream-Based Communication and Unix Sockets}
Unix Domain Sockets (\texttt{AF\_UNIX}, \texttt{SOCK\_STREAM}) provide full-duplex, reliable stream communication through the standard BSD socket API~\cite{stevens2004unix}. Unlike network sockets that traverse the TCP/IP network stack, Unix domain sockets transfer data entirely within kernel memory via socket buffers (\texttt{struct sk\_buff}). Unix domain sockets offer advanced capabilities, including file descriptor passing (\texttt{SCM\_RIGHTS}) and dynamic socket buffer auto-tuning. However, the internal socket layer introduces protocol overhead, state machines, and buffer allocation routines that can elevate per-message latency for small payloads.

\subsection{Existing IPC Performance Evaluations and Research Gap}
The primary baseline for this study is the empirical evaluation conducted by Danyliuk (2026)~\cite{danyliuk2026classical}, which analyzed round-trip latency across classical Linux IPC mechanisms in High-Performance Computing (HPC) scenarios. Danyliuk demonstrated that POSIX Shared Memory achieved the lowest latency for small payloads, while Unix Domain Sockets exhibited robust scaling for larger transfers.
However, existing studies exhibit critical limitations:
(1)~\textit{Lack of Virtualization Control:} Published benchmarks frequently conflate measurements gathered under virtualized hypervisors (WSL2, KVM) with bare-metal Linux results, masking underlying mechanism differences~\cite{renault2018efficient};
(2)~\textit{Absence of Tail-Latency Characterization:} Studies predominantly report mean latencies, obscuring worst-case scheduling jitter and high percentiles ($P_{95}, P_{99}$) critical for real-time systems; and
(3)~\textit{Static Evaluation Paradigms:} Prior work evaluates static transport bindings, offering no dynamic runtime mechanisms to exploit complementary transport strengths.
Mini-A-IPC directly bridges this gap by unifying six classical transports, an in-kernel driver mailbox, and an adaptive telemetry selector under rigorous experimental controls.
